\begin{trapbox}

	\begin{description}[style=nextline, leftmargin=2.8em, labelsep=0.5em, itemsep=0.35em]
		\item[\textbf{\textcolor{CTrap}{易错点一（定义域遗漏）}}]
			求正切函数定义域时，易漏掉分母不为零条件，或错误写成 $x\neq\frac{\pi}{2}+2k\pi$。
		\item[\textbf{\textcolor{CTrap}{正确理解（表达式正确）：}}]
			正切定义域为 $\{x\mid x\neq\frac{\pi}{2}+k\pi,k\in\mathbb{Z}\}$，排除点间隔 $\pi$。
		\item[\textbf{\textcolor{CTrap}{错因分析（周期混淆）：}}]
			误以为与正弦、余弦周期一致（$2\pi$），实际上正切周期为$\pi$，故排除点也以$\pi$为步长。
	\end{description}

	\medskip
	\begin{description}[style=nextline, leftmargin=2.8em, labelsep=0.5em, itemsep=0.35em]
		\item[\textbf{\textcolor{CTrap}{易错点二（周期缺绝对值）}}]
			求 $y=\sin(-2x)$ 的周期时，直接取 $T=\frac{2\pi}{-2}$，得到 $T=-\pi$，错误，周期应为正。
		\item[\textbf{\textcolor{CTrap}{正确理解（公式带绝对值）：}}]
			周期公式 $T=\frac{2\pi}{|\omega|}$，这里 $|\omega|=2$，故 $T=\frac{2\pi}{2}$，即 $T=\pi$。
		\item[\textbf{\textcolor{CTrap}{错因分析（忽视绝对值）：}}]
			误认为 $\omega$ 就是系数本身，未意识到周期必须为正数。
	\end{description}

	\medskip
	\begin{description}[style=nextline, leftmargin=2.8em, labelsep=0.5em, itemsep=0.35em]
		\item[\textbf{\textcolor{CTrap}{易错点三（单调区间混淆）}}]
			误以为余弦函数在 $[0,\pi]$ 上是增函数，实际是减函数。
		\item[\textbf{\textcolor{CTrap}{正确理解（余弦增区间）：}}]
			余弦函数增区间为 $[-\pi+2k\pi,2k\pi]$，减区间为 $[2k\pi,\pi+2k\pi]$。
		\item[\textbf{\textcolor{CTrap}{错因分析（记忆混淆）：}}]
			对余弦曲线图像记忆不清，或与正弦单调区间混淆。
	\end{description}

\end{trapbox}
